\chapter*{Til lesaren}
\addcontentsline{toc}{chapter}{Til lesaren}
\markboth{Til lesaren}{}

\begin{center}\itshape
Dette er ei historiebok. Men det er ei historie om eit fag som aldri\\ fann grunnen sin, fortald slik at du sjølv kan sjå kvar grunnen mangla.
\end{center}
\vspace{8mm}

Dei fleste designhistorier er sigerssoger. Dei fortel om rørsler som avløyste kvarandre, om meistrar som braut med forgjengarane sine, om stilar som modnast og institusjonar som voks; og dei fortel det som framsteg, som ein disiplin på veg mot seg sjølv. Denne boka fortel den same kronologien, frå laugshandverket gjennom den industrielle revolusjonen, Arts and Crafts, Werkbund, Bauhaus, Ulm, den amerikanske stylingen, den skandinaviske myten, metoderørsla, den semantiske vendinga, design thinking, og fram til overvakingskapitalismen, men ho les henne mot håra. Der sigerssoga ser ein disiplin som finn forma si, ser denne boka ein disiplin som gong på gong lova seg sjølv eit fundament og kvar gong leverte noko anna: ein stil, ein metode, ein pedagogikk, ein prosess, eit varemerke. Aldri grunnen sjølv.

Påstanden boka fører i bevis er enkel å seie og ubehageleg å høyre: \emph{designfaget har aldri nådd eit haldbart fagleg fundament}, og det har innretta seg slik at fråvêret aldri treng å bli innrømt. Eg kjem ikkje til å banke deg i hovudet med påstanden. Eg kjem til å fortelje historia, og lar deg lese diagnosen ut av to ting historia heile tida legg fram av seg sjølv: \emph{fråværa} og \emph{skadane}.

\section{Å lese eit fag gjennom det som manglar}

Eit modent fag kjennest att på kva det har. Medisinen har eit internasjonalt etisk kodeverk, eit organ som kan frata ein lege retten til å praktisere, ein delt metode, eit fagspråk presist nok til at ein lege i Oslo og ein i Tokyo meiner det same med dei same orda. Spør kva designfaget har av det same, og du møter ei rekkje fråvær. Det finst ikkje noko internasjonalt etisk råd for design med makt til å ekskludere. Det finst ingen felles designetikk utleidd av faget sjølv, berre etikk lånt utanfrå og bolta på som eit emne. Det finst ingen delt metode; det finst skular som er usamde om sjølve spørsmålet om design \emph{har} ein metode. Det finst ikkje noko felles fagspråk; orda faget feller dommar med, som «balanse», «integritet» og «det fungerer», tyder ulike ting i kvart atelier. Og det finst inga stemme: når verdsproblem som krev formgjeving står på dagsordenen, er det ikkje designfaget som blir kalla inn som ekspertise, slik klimavitskapen eller folkehelsa blir det. Faget har inga adresse å tale frå.

Desse fråværa er ikkje tilfeldige hol som tida vil fylle. Dei er \emph{systematiske}, og dei heng saman: kvart av dei følgjer av det same manglande fundamentet. Du kan ikkje ha eit etisk råd som ekskluderer for brot utan ein standard å bli ekskludert frå. Du kan ikkje ha ein delt metode utan eit delt omgrep om kva forma er og kva som verkar på henne. Du kan ikkje ha ei stemme utan ein kunnskap berre du eig. Difor er metoden i denne boka å lese fråværa som symptom. Kvar gong historia når eit punkt der eit modent fag ville hatt noko, og designfaget ikkje har det, noterer vi fråvêret og spør kvifor. Svaret er kvar gong det same, og at det er det same er sjølve poenget.

\section{Å lese eit fag gjennom det det skadar}

Det andre lesegrepet er hardare. Eit fag utan eit fundament som tvingar fram fleire omsyn enn det eine kunden bad om, byggjer det skadelege like villig som det gode, fordi det manglar språket til å skilje dei. Difor er skadane ikkje ein parentes i designhistoria; dei er prov. Boka følgjer ei line av skade gjennom heile kronologien (den planlagde foreldinga som amerikansk industridesign gjorde til forretningsmodell, eingongskulturen, mengda gjenstandar me eig og kastar), og ho endar der lina er tydelegast i dag: i selskapa som er mellom dei høgast verdsette i verda, og som har som uttalt mål å ekstrahere privat informasjon, fange merksemd og byggje avhengigheit. Det er ikkje misbruk av design. Det er design utført kompetent, langs den eine aksen oppdraget spesifiserte, av folk med designgrader, premiert på designkonferansar. At faget ikkje har noko internt språk for å kalle dette noko anna enn vellukka, er den skarpaste mogelege målinga av kva fundamentet ville gjort, og ikkje gjer.

Eitt tilfelle får eit eige kapittel, fordi det viser fråvêret innanfrå akademia sjølv: \emph{design thinking}. Her produserte faget endeleg noko det kunne eksportere, og eksporterte det med stor suksess, til økonomifaget, til handelshøgskulane, til konsulentbransjen. Men sjå kva som migrerte. Ikkje ein designfagleg kunnskap, for det fanst ingen å sende med; det som migrerte var ein tom prosedyre som kven som helst kunne ta i bruk fordi han ikkje kravde fagkunnskap. Og poda slo rot i økonomifaget, ikkje fordi økonomien hadde noko å tilby designen, men fordi designen ikkje hadde nokon grunn til å bli verande heime. Eit fag med eit fundament held på avkomet sitt. Eit fag utan eit, ser sitt einaste vellukka bidrag drive bort og feste seg der det er nokon ramme å feste det i.

\section{Korleis boka er bygd}

Boka har sju delar og tjuesju kapittel. \textsc{Del i} etablerer føresetnadene før ordet «design» fanst: handverket, den industrielle revolusjonen, og det fyrste forsøket på reform. \textsc{Del ii} til \textsc{iv} følgjer dei store rørslene og institusjonane som kvar lova eit grunnlag (Arts and Crafts, Werkbund, Bauhaus, Ulm, metoderørsla, Simon sin draum om ein designvitskap) og viser kvar dei stogga. \textsc{Del v} følgjer migrasjonen ut av faget: den semantiske vendinga, skjermen, og design thinking si ferd inn i økonomien. \textsc{Del vi} samlar skadane og fråværa i si tydelegaste form. \textsc{Del vii} gjer opp rekneskapen: plasserer faget i vitskapshistoria, og peikar, utan å love at faget vil gå dit, mot det einaste fundamentet eg veit om som kunne fylt holet, slik det er lagt fram i traktaten \textsc{Formlære}.

Kvart kapittel kan lesast for seg, og kvart kapittel endar med ein bolk \emph{Vidare lesing} for den som vil grave djupare i nett den perioden eller debatten. Saman er kapitla meinte å lesast som det dei er: ein samanhengande argumentasjon i historisk form. Du treng ikkje tru på konklusjonen for å ha nytte av historia. Men når du har lese henne, vil du ikkje lenger kunne sjå designhistoria som ei rein sigerssoge. Du vil sjå fråværa. Og fråværa er det mest ærlege portrettet faget har.

\vidarelesing{%
For oversyn over disiplinen og historiografien, sjå \textcite{margolin2015}, \textcite{forty1986}, \textcite{heskett1980} og \textcite{woodham1997}. For metode og sjølvforståing i designhistoria, \textcite{fallan2010} og \textcite{margolin1992histories}. Det teoretiske motstykket til denne boka er traktaten \textsc{Formlære} (i same prosjekt), som legg fram fundamentet boka her påviser fråvêret av.}
